\section{Lean architecture, repository reuse and completion path}\label{sec:lean}

\subsection{The target is a derivation in iPRA}
The relevant output type is
\[
            \Gamma\vdash[\mathsf{PRAAx}]\varphi,
       \qquad \varphi:\mathsf{Fm}\;\mathsf{PTm}\;n,
\]
not merely a proposition proved in Lean's ambient logic. The supplied
source files use the actual ProveIt notation for this judgment.
Here \texttt{PTm} is scoped term syntax, \texttt{Fm} is scoped
first-order formula syntax, and \texttt{Prf} is an inductive
intuitionistic natural-deduction derivation. \texttt{PRAAx} contains
exactly the arithmetic axioms of Section~\ref{sec:foundation}.

The existing syntax is arity-indexed. In schematic Lean notation:
\begin{verbatim}
inductive PR : Nat -> Type
  | zero : PR 0
  | succ : PR 1
  | proj {k} (i : Fin k) : PR k
  | comp {m k} (f : PR m) (gs : Fin m -> PR k) : PR k
  | prec {k} (f : PR k) (g : PR (k + 2)) : PR (k + 1)

inductive PTm (n : Nat) : Type
  | var (i : Fin n) : PTm n
  | app {k} (f : PR k) (ts : Fin k -> PTm n) : PTm n
\end{verbatim}
The finite functions in these constructors are tuples of syntax in the
metatheory. They are not arbitrary functions admitted into iPRA.
Formula constructors are equality of two terms, false, conjunction,
disjunction, implication, and the two quantifiers; the latter reduce
the number of free variables by one. Substitution shifts indices under
binders. \texttt{Fm.QF} recursively excludes both quantifier constructors.
The induction constructor requires an explicit proof of this property.
The appendix source closure includes these definitions and their
substitution lemmas, rather than replacing them with unspecified axioms.

\subsection{Essential data records}
The definition \texttt{PresentedCategory.Pres} uses the following exact
field types. Its associated \texttt{Laws} record contains derivation
constructors, not assertions that hold only in a semantic interpretation.
\begin{center}
\begin{tabular}{lll}
\toprule
Field & Syntax type & Intended arguments\\
\midrule
\texttt{Obj} & \texttt{Fm PTm 1} & $a$\\
\texttt{Hom} & \texttt{Fm PTm 3} & $a,b,f$\\
\texttt{Eqv} & \texttt{Fm PTm 4} & $a,b,f,g$\\
\texttt{idt} & \texttt{PTm 1} & $a$\\
\texttt{cmp} & \texttt{PTm 5} & $a,b,c,g,f$\\
\bottomrule
\end{tabular}
\end{center}
For example, its associativity field maps derivations typing $f,g,h$
to a derivation of the appropriate instantiated \texttt{Eqv} formula.
The human proof's endpoint guards are retained explicitly when needed;
the upstream minimal record does not supply every possible endpoint
closure rule automatically.

\texttt{Adjunction.Fun} supplies object and arrow terms of arities one
and three. A general natural transformation has a component term of
arity one and object-language typing and naturality laws. A family with
additional numerical parameters adds those parameters to the term arity.
\texttt{Yoneda.Fam} has the formulas \texttt{Par}, \texttt{Car},
\texttt{Sim} of arities one, three and four, and \texttt{act : PTm 5}.
Its law record supplies the presheaf action and congruence laws used in
the Yoneda calculation. For subsequent constructions requiring all
setoid equivalence or carrier-closure laws, those laws must also be
supplied; a minimal upstream record is not silently strengthened.

For a fixed diagram with objects $D_0,\ldots,D_{r-1}$ and arrows $D(u)$,
the cone formula at vertex $x$ and a coded component tuple $c$ is the
finite conjunction
\[
 \operatorname{Cone}_D(x,c)\equiv
 \bigwedge_i H_C(x,D_i,c_i)\ \wedge\
 \bigwedge_{u:i\to j}E_C(x,D_j,D(u)c_i,c_j).
\]
A selected limit record contains a term $L$, projection terms $p_i$,
a mediator term $m(x,c)$, and derivations of: the cone laws; typing of
$m$ under the cone premise; $p_i m=c_i$; and
\[
 H_C(x,L,v)\wedge\bigwedge_i E_C(x,D_i,p_i v,c_i)
           \ \longrightarrow\ E_C(x,L,v,m(x,c)).
\]
Each fixed conjunction is ordinary formula syntax. Products, equalizers
and colimits are specializations or duals. Varying diagram length is
handled by the finite-table Boolean representation, not by this fixed
conjunction as though it were a PR truth evaluator for arbitrary formulas.

A monad record adds fixed terms for $T$, $\eta$ and $\mu$, typing and
naturality derivations, and the three monad equations. The Kleisli
presentation has hom formula $H_C(a,Tb,f)$ and composition term
$\mu_cT(g)f$. The Eilenberg--Moore object formula on $(a,h)$ is
\[
 O_C(a)\wedge H_C(Ta,a,h)\wedge h\eta_a=\id_a
                         \wedge h\mu_a=hT(h),
\]
with its displayed hom predicate from Section~\ref{sec:monads}.
The Karoubi object formula uses $e^2=e$, and its identity term returns
$e$. These arithmetical formulas may be undecidable, but their identity
and composition constructors are PR. To complete their Lean proofs,
one expands the finite calculations in the report into the already
available equality and substitution derivation constructors.

\subsection{What was actually inspected in ProveIt}
All source reuse is pinned to commit
\texttt{72d37f6d06fa7b5bd0a46732060446d1cffa3f2d}, retrieved on
20 November 2026 through the GitHub connector. The repository specifies
Lean \texttt{v4.32.0} and Mathlib \texttt{v4.32.0}.
The package contains 104 upstream Lean files forming the complete
in-repository import closure of the selected categorical modules, their
license, and a source manifest. Mathlib and the Lean executable are
external build dependencies.

\begin{longtable}{@{}>{\raggedright\arraybackslash}p{0.30\linewidth}>{\raggedright\arraybackslash}p{0.63\linewidth}@{}}
\toprule
Module & Inspected object-language content and limitations\\
\midrule
\endhead
\texttt{PresentedCategory} & Category formulas and PR terms; inverse
uniqueness, sections and retractions, terminal uniqueness; a concrete
additive-monoid instance.\\
\texttt{Yoneda} & Presheaf and transformation law records;
\texttt{yoneda}, \texttt{yoneda\_closed}, \texttt{repr\_natural},
\texttt{repr\_eval}, \texttt{rep\_laws}, representable specialization.\\
\texttt{Adjunction} & Functor and adjunction records; both hom-map
round trips; construction of the unit and counit from a natural
correspondence, with naturality and triangles.\\
\texttt{SplitEquivalence} & Explicit hom lifting, object representatives,
quasi-inverse, inverse equations and naturality.\\
\texttt{FiniteLimits} & Rank/select construction at a fixed diagram
shape; assumes a retention test with soundness and completeness and
positive size conditions. The general table-built test is not supplied.\\
\texttt{FiniteSetClosed} & Evaluation, currying and their round trips
under positive-cardinality hypotheses. Empty cases need extension.\\
\texttt{FiniteCoequalizer} & Sequential label merging and its
factorization calculations in the source's stated finite-set encoding.\\
\texttt{KanFormula} & Representative-level left and right Kan formulas,
their compatibility and naturality; not an unrestricted existence
theorem for arbitrary comma (co)limits.\\
\texttt{SetoidInterpretation} & Domain- and equality-respecting
translation of object-language derivations.\\
\texttt{UniversalEvaluator}\linebreak\texttt{Obstruction} & An external no-universal-PR-
evaluator theorem and, separately, an actual iPRA diagonal sentence.\\
\texttt{FiniteCoverSkeleton}, \texttt{TwoClassObstruction} & Positive
bounded representative selection and the precise obstruction to
extracting it merely from a bound on equivalence classes.\\
\bottomrule
\end{longtable}

The repository describes its \texttt{Derivations} library as kernel-checked.
This packet has inspected those pinned sources but has \emph{not} rebuilt
them locally. A proof-shaped source declaration is not by itself a fresh
build receipt. Nor should an older research packet's description of
``uncompiled sketches'' be applied to a newer \texttt{Derivations} module
without reading the latter. These distinctions matter when reporting reuse
\cite{repo-readme,repo-presented,repo-yoneda,repo-adjunction}.

\subsection{New Lean adapters and a useful interface refinement}
The source \texttt{Yoneda.NatLaws} asks one term
$\alpha(p,A,x,f)$ to obey its laws uniformly for every family parameter
$p$ and representing object $A$. This is a valid sufficient interface,
but stronger than the data of one transformation $h_A\Rightarrow P_p$
at fixed $p,A$. It should not be mistaken for that local interface.

The new \texttt{CategoryIPRA/LocalYoneda.lean} defines
\texttt{AtLaws} at fixed $p,A$ and object-language context $\Gamma$.
It repeats the short Yoneda derivation with these local premises and
provides conversion from the upstream uniform record.
\texttt{Representability.lean} derives faithfulness and the two inverse
equations from a specified natural isomorphism, using this localized
record. \texttt{Density.lean} derives value typing, canonical-value
evaluation, generator round trip and reindexing for the density
construction. Their conclusions are actual iPRA derivation judgments;
they introduce no universal evaluator and no new axiom declaration.

\textbf{Verification status.} All three new Lean modules are uncompiled
integration drafts. No Lean executable was available in the generation
environment. They contain no \texttt{sorry} or \texttt{admit}, but their
elaboration and kernel checking remain to be performed. The completed
human proofs are separate evidence; no statement in this report treats
the drafts as a completed mechanization. The source-structure audit
parsed the Lake configuration and found no missing in-repository modules
among 192 imports across 110 Lean source files. External import names lie
under the specified Mathlib dependency; its installed module files were
not checked because the build environment was unavailable.

\subsection{A concrete path to completion}
The archive's \texttt{lean/formalization\_roadmap.md} gives record fields,
dependencies and coding obligations in more detail. The following order
keeps each step mathematically reviewable.
\begin{enumerate}
\item Build the pinned source closure and the three new adapters; resolve
elaboration issues and record the actual kernel and axiom-audit output.
Close the desired local derivations by the existing universal-closure
machinery, constructing law records under the required contexts.
\item Add generic natural transformations, selected cones and universal
operations as syntax records. Port their finite equational proofs using
the supplied category and functor law constructors.
\item Complete density's tuple coding and cocone record. The difficult
logical point is already removed by direct equality of values; no
unbounded generated-equivalence closure is needed.
\item Port the base-$(m+1)$ finite tables, the Boolean diagram tester,
rank/select, and all zero-cardinality cases. State every loop invariant
as a concrete quantifier-free formula. The executable model provides
reference outputs, not proofs of these invariants.
\item Encode the monad, Kleisli, algebra and idempotent formulas and
operations exactly as above. Each law is a finite equation chain;
algebra limits use the selected limiting mediator and its uniqueness.
\item Encode the parameterized comma diagrams and their selected
(co)limits. Define the Kan functor and mediator by the formulas in
Section~\ref{sec:kan}. Preserve the explicit existence assumptions.
\item Audit every theorem as either a closed iPRA derivation, a derivation
under named hypotheses, or an external schema. In particular, do not
replace code-level PR certificates by unrestricted semantic Lean
functions, and do not infer an internal soundness principle from a
metatheoretic proof transformation.
\end{enumerate}

With the pinned toolchain and Mathlib available, the intended commands are
\begin{verbatim}
cd lean
lake update
lake exe cache get
lake build CategoryIPRA
lake env lean CategoryIPRA/Audit.lean
\end{verbatim}
These are reproduction instructions. They have not been presented as
successful executions in this environment.

\subsection{Executable validation and its limits}
The exact-arithmetic model in \texttt{code/finite\_models.py} uses integer
tables, including empty tables and base-$(m+1)$ encodings. An actual run
passed 22,016 exponential beta cases and 22,016 eta cases; 1,544 product
beta cases and 1,544 eta cases; 60 valid and 84 invalid table-code cases;
and eight limit/colimit example diagrams, whose universal properties were
exhaustively checked against test objects of cardinality at most three.
The diagrams include the empty diagram, empty objects, an empty equalizer,
a pullback and a quotient requiring transitive identifications.

The powerset presheaf on the finite full subcategory with cardinalities
$0,1,2$ gives a more discriminating Yoneda test: all 1,024 possible
component families were checked, and precisely four were natural,
corresponding to the four elements of the powerset of $[2]$.
Two further examples use the group and zero monoid of order two.
The option monad's Kleisli operations passed 23 unit cases and 2,457
associativity cases.

The executable colimit implementation uses bounded Warshall closure;
the human proof gives sequential merging. Both realize the same
finite-path quotient. This independence is useful for catching example
errors, but tests of finitely many models do not prove an unrestricted
theorem or verify an iPRA proof checker. The JSON receipt records the
actual finite scope and all counts.
